% ===========================================================================
%  Apendice B — Dicionario de features
% ===========================================================================
\chapter{Dicionário completo de \textit{features}}
\label{ap:features}

O vetor de entrada do modelo tem 23 dimensões, na ordem exata abaixo. A
\textbf{ordem é significativa}: ela é a ordem esperada pelos artefatos
treinados. Reordenar a lista sem retreinar produz predições silenciosamente
erradas.

\section{Grandezas brutas (18)}

\begin{tabularx}{\linewidth}{@{}L{8mm}L{44mm}L{18mm}X@{}}
\toprule
\headrow \thd{\#} & \thd{Nome} & \thd{Unidade} & \thd{Significado físico}\\
\midrule
\endhead
1 & \cd{z\_rms\_velocity\_mm\_s} & \si{\milli\meter\per\second} &
  Velocidade eficaz de vibração no eixo Z (axial).\\
\zebra 2 & \cd{x\_rms\_velocity\_mm\_s} & \si{\milli\meter\per\second} &
  Velocidade eficaz no eixo X (radial). É o indicador normativo de severidade
  global~\cite{iso10816}.\\
3 & \cd{temperature\_c} & \si{\degreeCelsius} &
  Temperatura medida no ponto do sensor. \textit{Feature} mais influente do
  classificador --- e a mais suspeita de funcionar como carimbo de campanha
  (Seção~\ref{sec:vazamento}).\\
\zebra 4 & \cd{z\_peak\_acceleration\_g} & $g$ &
  Aceleração de pico no eixo Z. Sensível a impactos.\\
5 & \cd{x\_peak\_acceleration\_g} & $g$ &
  Aceleração de pico no eixo X.\\
\zebra 6 & \cd{z\_peak\_vel\_comp\_freq\_hz} & \si{\hertz} &
  Frequência da componente dominante de velocidade em Z.\\
7 & \cd{x\_peak\_vel\_comp\_freq\_hz} & \si{\hertz} &
  Idem, eixo X. Base da \textit{feature} derivada de ordem. \SI{61}{\hertz}
  responde por \SI{48.7}{\percent} dos registros --- distribuição quase
  degenerada.\\
\zebra 8 & \cd{z\_rms\_acceleration\_g} & $g$ &
  Aceleração eficaz em Z (banda completa).\\
9 & \cd{x\_rms\_acceleration\_g} & $g$ &
  Aceleração eficaz em X (banda completa).\\
\zebra 10 & \cd{z\_kurtosis} & --- &
  Curtose do sinal em Z. Teoricamente indicativa de
  impactos~\cite{randall2011tutorial}; neste conjunto de dados varia entre
  \num{2,65} e \num{2,75} em \emph{todas} as famílias.\\
11 & \cd{x\_kurtosis} & --- &
  Curtose em X, com a mesma limitação.\\
\zebra 12 & \cd{z\_crest\_factor} & --- &
  Razão pico/eficaz em Z. Varia entre \num{3,73} e \num{3,85} em todas as
  famílias.\\
13 & \cd{x\_crest\_factor} & --- &
  Fator de crista em X.\\
\zebra 14 & \cd{z\_peak\_velocity\_mm\_s} & \si{\milli\meter\per\second} &
  Velocidade de pico em Z.\\
15 & \cd{x\_peak\_velocity\_mm\_s} & \si{\milli\meter\per\second} &
  Velocidade de pico em X. Colinear com a velocidade eficaz em X --- razão $F$
  idêntica (\num{1193.6}), porque o fator de crista é quase constante.\\
\zebra 16 & \cd{z\_high\_freq\_rms\_accel\_g} & $g$ &
  Aceleração eficaz de alta frequência em Z.\\
17 & \cd{x\_high\_freq\_rms\_accel\_g} & $g$ &
  Idem, eixo X. Base da razão de alta/baixa frequência.\\
\zebra 18 & \cd{rpm} & rpm &
  Rotação do eixo. Variável de \emph{condição operacional}: estratificar por ela
  é obrigatório em qualquer leitura de severidade
  (Seção~\ref{sec:severidade}).\\
\bottomrule
\end{tabularx}

\section{\textit{Features} derivadas (5)}

\begin{tabularx}{\linewidth}{@{}L{8mm}L{40mm}L{40mm}X@{}}
\toprule
\headrow \thd{\#} & \thd{Nome} & \thd{Definição} & \thd{Motivação}\\
\midrule
19 & \cd{x\_peak\_order} &
  $\min\!\left(\dfrac{f_{\text{pico},X}}{\max(\text{rpm}/60,\,10^{-6})},\,50\right)$ &
  Ordem do pico no eixo X. Os manuais descrevem assinaturas em múltiplos da
  rotação, não em \si{\hertz} absoluto.\\
\zebra 20 & \cd{z\_peak\_order} &
  Idem, eixo Z &
  Ordem do pico no eixo axial.\\
21 & \cd{radial\_axial\_ratio} &
  $\dfrac{v_{\text{RMS},X}}{v_{\text{RMS},Z} + 10^{-6}}$ &
  Desbalanceamento é predominantemente radial; desalinhamento e rotor inclinado
  produzem componente axial apreciável.\\
\zebra 22 & \cd{hf\_lf\_ratio\_x} &
  $\dfrac{a_{\text{HF,RMS},X}}{a_{\text{RMS},X} + 10^{-6}}$ &
  Defeitos localizados em rolamentos excitam preferencialmente alta
  frequência.\\
23 & \cd{hf\_lf\_ratio\_z} &
  Idem, eixo Z &
  Mesma motivação, eixo axial.\\
\bottomrule
\end{tabularx}

\section{Colunas descartadas}

\begin{tabularx}{\linewidth}{@{}L{50mm}X@{}}
\toprule
\headrow \thd{Coluna descartada} & \thd{Motivo}\\
\midrule
Velocidades em \si{\inch\per\second} & Transformação linear exata das
  contrapartes em \si{\milli\meter\per\second} (fator \num{25,4}). Inclusão
  produziria colinearidade perfeita sem acrescentar informação.\\
\zebra Temperatura em \si{\degF} & Transformação afim exata da temperatura em
  \si{\degreeCelsius}.\\
\cd{fault} (rótulo bruto) & É o \emph{alvo}, não \textit{feature}. Preservado em
  \cd{events.raw\_fault} para auditoria e para identificar a campanha de coleta
  na análise de vazamento.\\
\zebra \cd{id} e \cd{created\_at} & Identificador e instante não são sintomas
  mecânicos. Incluí-los seria vazamento explícito, dado que o instante
  identifica a campanha.\\
\bottomrule
\end{tabularx}

\section{Poder discriminativo observado}

Razão entre a variância inter-famílias e a variância intra-família (estatística
no espírito do $F$ de análise de variância), calculada sobre as famílias com ao
menos 30 eventos.

\begin{tabularx}{\linewidth}{@{}XR{28mm}L{56mm}@{}}
\toprule
\headrow \thd{\textit{Feature}} & \thd{Razão $F$} & \thd{Leitura}\\
\midrule
velocidade RMS X & \num{1193.6} & Melhor discriminador do conjunto.\\
\zebra velocidade de pico X & \num{1193.6} & Colinear com a anterior.\\
frequência de pico X & \num{1106.8} & Alto valor, mas distribuição quase
  degenerada.\\
\zebra temperatura & \num{885.3} & Discrimina --- e também carimba a campanha.\\
rotação & \num{208.3} & Condição operacional, não sintoma.\\
\zebra curtose Z & \num{189.2} & \textbf{Baixo}: contraria a previsão dos
  manuais.\\
aceleração de pico Z & \num{188.2} & Baixo.\\
\bottomrule
\end{tabularx}
